\chapter{Design the keys}

The point of architecture here is not elegance. It is
that a coerced person, in a short window, cannot empty
the thing that matters. The victim should not have to
out-tough the attacker. The system should fail closed
without requiring heroics.

The figures below are examples from the source, not a
universal policy. Size them to your actual balances and
your actual attackers.

\section{Five principles}

\textbf{Minimization.} Cap what any one coerced person
can authorize in the time an attack usually lasts. Even
full compliance should yield a loss you can survive.
Make the attacker stay longer. Time is how they get
caught.

\textbf{Compartmentalization.} Split keys, identities,
devices, places, and roles. A burned hot wallet should
not point at cold storage. Signers should not all live
in the same city, or all report to the same person.

\textbf{Controlled surrender.} Give the attacker
something true and small they can take now, in a range
that matches what the public can already see. Keep the
real reserve behind delay and extra people, and make
that delay visible. A lie that collapses under stress is
worse than a small honest wallet.

\textbf{Resilience.} You will need to rebuild. No single
backup, no single restorer, no blame circus that leaks
the rest of the story.

\textbf{Prioritization.} Human safety beats assets. No
rule should reward physical resistance. Wallet policy
yields when someone is being hurt.

\section{Wallet tiers}

Think in tiers with different jobs, not one ``the
wallet.''

\textbf{Hot (example: 0--1\%).} Daily ops, gas, small
payments. Strict daily limits. No path from here into
cold.

\textbf{Warm (example: 1--5\%).} A few days of payroll or
vendors. Extra approval and delay as size grows.

\textbf{Cold (example: 5--25\%).} Hours to a few days,
hardware signers, stored off the home and the office.

\textbf{Deep-cold (the rest).} Strategic reserve.
Recovery that takes legal process or a ceremony measured
in days or weeks, not a PIN in a kitchen.

Those percentages are prompts. A solo trader and a DAO
treasury should not share a template.

\section{Multisig, MPC, geography}

For material movement, spread signing so that grabbing
one person does not complete the payment. The source
talks in ranges such as 3-of-5 or 4-of-7. The number
matters less than the independence: different places,
different risk, different ways of being reached.

MPC is the same idea with shares instead of distinct
keys. The requirement is identical. No single share
spends.

\textbf{Do not keep hardware wallets, seeds, or signer
devices at the primary residence.} Home is where wrench
attacks happen. Temporary hot access via a secure
element is a different problem from storing the reserve
under the bed.

Offsite means a vault, a bank box, or a trusted location
that is not on your daily loop. Deep-cold can be an
institutional custodian. Protocol treasuries can add
timelocks, pause, or an emergency path that a coerced
admin cannot skip.

\section{A wallet you can hand over}

Keep one wallet whose balance you can surrender without
lying, sized to look real against public activity. The
source's example range is a few thousand dollars, plus
a pre-funded exchange account around \$10k for personas
that would plausibly have one. Again: examples. A public
trader with obvious volume cannot hand over \$400 and
expect to be believed.

\textbf{Make the real reserve's delay obvious.}
Timelocks and extra signers the attacker can see. The
goal is to remove the incentive to keep hitting you for
the rest.

If the device supports a duress PIN, wire it to a silent
alert for a pre-agreed contact. Silent means no extra
call or text the attacker watches you send. This is the
control that gets people killed when it is amateurish.
Rehearse it with people who do crisis work, or do not
use it.

\section{Ceremonies and rotation}

Generate material in a place that is not anyone's house,
with more than one person, and without a whole seed
sitting on a table. Split shards. Record the ceremony if
your threat model wants an audit trail.

Review who can still sign. Refresh hardware on a
calendar, and after travel that raised the physical
risk. After an incident, assume a 24-hour replacement
path already exists. Inventing one while someone is in
hospital is how you reuse a compromised signer.
